\chapter{Родитель расщепления масс гиперзарядового проектора
\texorpdfstring{$R_2$}{R2}}
\label{ch:v10-h15-r2-hypercharge-mass-parent}

Предыдущий гейт дал точный проектор на
$R_2\oplus\overline R_2$, но не динамический член потенциала. Здесь
проверяется самый экономный mass-parent, использующий только уже известный
оператор $Q=6Y$. Его положительный спектральный gap равен
\begin{equation}
 \boxed{G_Y=49I-Q^2}
 =\operatorname{diag}(40,40,0,48,48,0,40,40).
\end{equation}
Он имеет ранг $6$, нуллитет $2$ и ядро ровно на
$R_2\oplus\overline R_2$. Поэтому для квадратичной части потенциала
естественна условная семья
\begin{equation}
 V_2(\Sigma)=\frac12\Sigma^\dagger
 \bigl(\kappa G_Y-\mu^2I\bigr)\Sigma,
 \qquad \kappa>0.
\end{equation}

Массы трёх спектральных классов $Q^2=49,9,1$ равны соответственно
\begin{equation}
 -\mu^2,\qquad 40\kappa-\mu^2,\qquad 48\kappa-\mu^2.
\end{equation}
Следовательно, только $R_2$-пара тахионна тогда и только тогда, когда
\begin{equation}
 \boxed{0<\frac{\mu^2}{\kappa}<40}.
\end{equation}
При положительной квартальной стабилизации вдоль $R_2$ это действительно
задаёт требуемое направление конденсации. Например, точный свидетель
$(\kappa,\mu^2)=(1,20)$ даёт диагональ
$(20,20,-20,28,28,-20,20,20)$ и сигнатуру $(2,0,6)$ в порядке
negative/zero/positive.

Оператор $G_Y$ совместим с Real-обменом сопряжённых секторов, но его
коммутатор с полным $SU(2)_R$-flip имеет ранг $4$. Это ожидаемо: gap
каноничен после выбора гиперзарядового SM-направления, а не до нарушения
Pati--Salam-симметрии.

Главный отрицательный результат относится к происхождению коэффициентов.
Текущий унаследованный gap-Hessian имеет ранг $0$, а карта происхождения
пары $(\kappa,\mu^2)$ нулевая. Условие $0<\mu^2/\kappa<40$ является открытым
интервалом, а не предсказанием отношения. Свидетель $20$ демонстрирует
непротиворечивость архитектуры, но не выбирается теорией. Абсолютный
масштаб также остаётся свободным.

\begin{theorem}[Условное гиперзарядовое расщепление]
После выбора SM-гиперзаряда оператор $G_Y=49I-(6Y)^2$ является
положительным спектральным gap с ядром ровно
$R_2\oplus\overline R_2$. Семейство
$\kappa G_Y-\mu^2I$ имеет ровно две отрицательные и шесть положительных
мод тогда и только тогда, когда $\kappa>0$ и
$0<\mu^2/\kappa<40$.
\end{theorem}

\begin{nogocriterion}[Интервал не выбирает родителя]
Нельзя считать mass-splitting физически выведенным, выбрав произвольную
точку внутри допустимого интервала. Требуются независимое происхождение
$\kappa$, тахионного отношения $\mu^2/\kappa$, положительного квартального
члена и абсолютного масштаба.
\end{nogocriterion}

\begin{protocol}\begin{enumerate}
\item спектр $G_Y$: \textbf{$0^2,40^4,48^2$};
\item ранг/нуллитет gap: \textbf{$6/2$};
\item Real-совместимость: \textbf{да};
\item полный $SU(2)_R$-дефект: \textbf{rank $4$};
\item окно $R_2$-only instability: \textbf{$0<\mu^2/\kappa<40$};
\item точный условный свидетель: \textbf{$\mu^2/\kappa=20$};
\item inherited gap-Hessian rank: \textbf{$0$};
\item происхождение $(\kappa,\mu^2)$: \textbf{$0/2$};
\item physical origin: \textbf{$1/3$};
\item следующий гейт: \textbf{аудит происхождения коэффициента гиперзарядового gap}.
\end{enumerate}\end{protocol}

Машинный сертификат сохранён в
\path{s2t_v10_particle_wrinkle_dislocation_callias_equal_charge_twist_m4_h15_r2_hypercharge_projector_mass_splitting_parent_origin_gate_results.json}.